\scanpage{3}
\mspage{1}
Soient $X$ et $Y$ deux surfaces topologiques \inserted{(connexes pour simplifier)} orientées ``de type fini'' (i.e. déduites de surfaces compactes en enlevant un nb fini de points). On considère
\[
\begin{aligned}
(1)\qquad \operatorname{Hom}_{\mathrm{isot}}(X,Y)
&=\text{ensemble des composantes connexes \inserted{par arcs}}\\
&\quad\text{de l'espace des applications finies}\\
&\quad\text{étales de $X$ dans $Y$ \deleted{respectant}}\\
&\quad\deleted{\text{l'orientation}}
\end{aligned}
\]
Si $f$ et $g$ sont deux applications finies étales de $X$ dans $Y$, alors $f$ et $g$ \inserted{sont isotopes i.e.} définissent la $=$ \deleted{\unsure{applicat|appli}} application de $X$ en $Y$ dans la cat. isotopique, ssi nous pouvons trouver une application continue \unsure{fin étal|fin ét}
\deleted{$F$ de $X\times I$}
\[
(2)\qquad F:X\times I\longrightarrow Y\times I\qquad I=[0,1]
\]
commutant aux projections sur $I$, et telle que
\[
(3)\qquad F(x,0)=(f(x),0)\qquad F(x,1)=(g(x),1)\qquad\forall x\in X
\]
Comme le foncteur
\[
Y'\longmapsto Y'\times I
\]
des rev. finis étales de $Y$ vers les rev. finis étales de $Y\times I$ est pl. fidèle, on voit que \deleted{\unsure{il exi}} pour $F$ donné, il doit exister un rev. fini étale $Y'$ de $Y$, unique à isom. unique près, et une \unsure{homéom.}
\[
(4)\qquad \alpha:X\times I\xrightarrow{\sim}Y'\times I
\]
rendant commutatif le diag. \unsure{comm.|suiv.}
\diagram{
\begin{tikzcd}[ampersand replacement=\&]
(5)\quad X\times I \arrow[rr,"\alpha", "\sim"'] \arrow[dr,"F"'] \&\& Y'\times I \arrow[dl,"p\times\operatorname{id}_I"]\\
\&Y\times I\&
\end{tikzcd}
}
\margin{où $p:Y'\mathbin{-}Y$ la projection}
\margin{2}
